\documentclass[12pt]{article}
\usepackage[margin=1in]{geometry}
\usepackage{booktabs,threeparttable}

\begin{document}

\providecommand{\StructuralRobustnessPath}{appendix/structural-robustness}

\subsection{Alternative Measures of Observability}
\label{app:alternative-observability}

The baseline model measures observability using the share of recognized peers
for whom respondents give a definite Yes or No answer about deworming. We
re-estimate the complete structural model using two alternative inputs. The
first uses the share of recognized peers whose status is classified correctly
when linked to administrative take-up records. The second uses perceived
second-order observability: respondents' beliefs about whether other community
members know who dewormed. We also make the treatment of missing first-order
belief responses explicit by retaining those respondents and marginalizing
their unobserved belief outcomes.

Table~\ref{tab:alternative-observability-multipliers} shows that the multiplier
pattern is stable. Under correct classification, the multiplier remains above
one for Control and Calendar and is generally below one for Ink; the Bracelet
point estimates are close to one and fall with distance. Under perceived
second-order observability, the estimates closely track the baseline, with
amplification for Control and Calendar and mitigation for Bracelet and Ink at
longer distances.

In the present observability model, explicit marginalization of missing belief
outcomes is algebraically equivalent to the baseline target. Missing rows make
no likelihood contribution, and predicted observability depends only on
treatment arm and community distance. Hence all respondents in a given
community have the same predicted probability, so averaging predictions over
observed and missing respondents produces the same community observability
schedule. Panel D records this equality rather than presenting a poorly mixed
duplicate fit as independent evidence.

\begin{table}[htbp]
  \centering
  \small
  \caption{Social multipliers under alternative measures of observability}
  \label{tab:alternative-observability-multipliers}
  \begin{threeparttable}
    \begin{tabular}{lccc}
\toprule
Treatment & 500m & 1500m & 2500m \\
\midrule
\multicolumn{4}{l}{\textit{Panel A: First-order definite classification (baseline)}} \\
Bracelet & 0.89 & 0.83 & 0.78 \\
 & (0.68, 1.04) & (0.68, 0.96) & (0.67, 0.90) \\
Calendar & 1.25 & 1.22 & 1.21 \\
 & (1.12, 1.48) & (1.05, 1.47) & (1.00, 1.53) \\
Ink & 0.92 & 0.86 & 0.81 \\
 & (0.80, 1.03) & (0.75, 0.98) & (0.70, 0.95) \\
Control & 1.43 & 1.48 & 1.56 \\
 & (1.25, 1.75) & (1.25, 1.88) & (1.27, 2.03) \\
\addlinespace
\multicolumn{4}{l}{\textit{Panel B: Correct classification}} \\
Bracelet & 1.01 & 0.95 & 0.90 \\
 & (0.76, 1.36) & (0.73, 1.25) & (0.69, 1.19) \\
Calendar & 1.45 & 1.35 & 1.30 \\
 & (1.14, 2.04) & (1.08, 1.83) & (1.04, 1.73) \\
Ink & 0.95 & 0.90 & 0.86 \\
 & (0.77, 1.16) & (0.73, 1.11) & (0.68, 1.08) \\
Control & 1.59 & 1.49 & 1.42 \\
 & (1.25, 2.28) & (1.19, 2.10) & (1.15, 1.97) \\
\addlinespace
\multicolumn{4}{l}{\textit{Panel C: Perceived observability (second order)}} \\
Bracelet & 0.89 & 0.83 & 0.78 \\
 & (0.61, 1.09) & (0.64, 0.97) & (0.65, 0.90) \\
Calendar & 1.34 & 1.27 & 1.25 \\
 & (1.13, 1.69) & (1.08, 1.60) & (1.03, 1.61) \\
Ink & 0.98 & 0.90 & 0.84 \\
 & (0.83, 1.11) & (0.76, 1.02) & (0.70, 0.97) \\
Control & 1.48 & 1.48 & 1.53 \\
 & (1.22, 1.88) & (1.21, 1.88) & (1.21, 1.97) \\
\addlinespace
\multicolumn{4}{l}{\textit{Panel D: First-order observability, missing data marginalized}} \\
Bracelet & 0.89 & 0.83 & 0.78 \\
 & (0.68, 1.04) & (0.68, 0.96) & (0.67, 0.90) \\
Calendar & 1.25 & 1.22 & 1.21 \\
 & (1.12, 1.48) & (1.05, 1.47) & (1.00, 1.53) \\
Ink & 0.92 & 0.86 & 0.81 \\
 & (0.80, 1.03) & (0.75, 0.98) & (0.70, 0.95) \\
Control & 1.43 & 1.48 & 1.56 \\
 & (1.25, 1.75) & (1.25, 1.88) & (1.27, 2.03) \\
\bottomrule
\end{tabular}

    \begin{tablenotes}\footnotesize
      \item \textit{Notes:} Entries are posterior medians of the social
      multiplier at the indicated distance. Parentheses contain 95 percent
      posterior credible intervals. Panels A--C use four-chain posterior fits.
      Panel D is identical to Panel A by construction under the current
      observability model, as described in the text.
    \end{tablenotes}
  \end{threeparttable}
\end{table}

\end{document}
